\documentclass{article}
\usepackage{amsmath, amssymb, amsthm}
\usepackage{hyperref}
\usepackage{geometry}
\geometry{margin=1in}

\title{Erd\H{o}s Problem \#1004}
\date{Last updated: 2025-09-07}
\author{Source: erdosproblems.com}

\begin{document}
\maketitle

\noindent\textbf{Status:} OPEN \quad \textbf{No prize}

\noindent\textbf{Tags:} number theory

\medskip
\noindent\textbf{Problem Statement:}

Let $c>0$. If $x$ is sufficiently large then does there exist $n\leq x$ such that the values of $\phi(n+k)$ are all distinct for $1\leq k\leq (\log x)^c$, where $\phi$ is the Euler totient function?

\bigskip
\noindent\textbf{Remarks:}

Erd\H{o}s, Pomerance, and S\'{a}rk\"{o}zy \cite{EPS87} proved that if $\phi(n+k)$ are all distinct for $1\leq k\leq K$ then\[K \leq \frac{n}{\exp(c(\log n)^{1/3})}\]for some constant $c>0$. 

See [945] for the analogous problem with the divisor function.

\bigskip
\noindent\textbf{References:}

[EPS87] Erd\H{o}s, Paul and Pomerance, Carl and S\'{a}rk\"ozy, Andr\'{a}s, On locally repeated values of certain arithmetic functions.
{III}. Proc. Amer. Math. Soc. (1987), 1--7.

\bigskip
\noindent\small{Source: \url{https://www.erdosproblems.com/1004}}
\end{document}
